\section{Selection, Physical Feasibility, and Three-Tier State}
\label{app:selection}
\subsection{Blockization and invariants}
A \texttt{KVBlockDescriptor} records layer, KV head, modality, physical source
indices, original logical positions, token IDs when available, image/frame/page/
region metadata, K/V dtype and device, payload bytes, and a mandatory flag.
Fixed-size token blockization never mixes layers or KV heads. A tail or
modality boundary can produce a shorter block, so actual cost is recorded
rather than inferred from the configured size alone.

Every source position in each layer/head belongs to exactly one source block.
Logical positions are monotone within a block. Mandatory source blocks cannot
be removed by a later feasibility fill, prototype reassignment, or repair
eviction. The membership rules require care: exact members and prototype
source members are disjoint, while residuals intentionally duplicate the
original members represented by prototypes. Calling this intentional overlap
a membership conflict would make recovery impossible. Exact-only eviction can
leave unrepresented source blocks, which are recorded as discarded rather
than silently counted as residuals.

\subsection{Utility components and signed weights}
Within one layer/head stratum, each nonnegative signal is independently
normalized to unit sum when its sum is nonzero. Attention contribution is
estimated from forecast mass; a simple value contribution is
$c_i^v=a_i\lVert\bar v_i\rVert_2$. Value novelty is
\begin{equation}
 n_i=1-\max_{j\in\mathcal N(i)}\max\{0,\cos(\bar v_i,\bar v_j)\}.
\end{equation}
An empty neighborhood requires a declared convention. Raw terms are retained
so the normalized score can be audited. Centrality, modality rarity, redundancy,
and mandatory priority are additional diagnostics or configured policy inputs;
they must not be claimed as extra additive terms if the selected objective
does not actually use them.

The implementation uses literal minus signs for the value and novelty terms
and for coverage. Negative coefficients turn these terms into rewards.
For explanatory purposes one may set $\alpha_v=-\lambda_v$ and
$\alpha_g=-\lambda_g$ and discuss nonnegative reward coefficients, but the
manifest must retain the actual resolved values. Changing sign in prose
without changing the configuration would describe a different method.

\subsection{Greedy selection and its limits}
Mandatory blocks are selected in deterministic order before optional choices.
For feasible candidate $i$, the marginal gain is
$\Delta_i(S)=F(S\cup\{i\})-F(S)$. In a byte budget, a greedy priority can use
$\Delta_i/c_i$. Ties are resolved by raw gain, lower cost, then stable block ID.
Cached marginal bounds are valid only for the relevant diminishing-returns
objective. Positive coverage penalties violate the intended lazy-greedy
condition and must not be silently accepted.

The unconstrained-cardinality approximation result for normalized monotone
submodular maximization \citep{nemhauser1978} does not establish optimality for
mandatory, per-head, padded, modality-constrained byte packing. Tiny synthetic
graphs can be exhaustively enumerated to report an objective gap and test
diminishing returns. This is an algorithm test, not a benchmark score.

The production execution variant may use vectorized per-route top-k for
single-token equal-cost blocks. It preserves deterministic ordering and avoids
general per-head dynamic programming. It is not automatically an exact
optimizer of graph coverage. The trace must identify which selector ran, which
heads entered it, the selected marginal/reason for each block, and any fallback.

\subsection{Why the largest feasible cache is not simply the cap}
For equal-width dense heads with lengths $n_{\ell h}$, a layer may allocate
storage according to $\max_h n_{\ell h}$. The selected-source sum can then be
smaller than actual storage. In a uniform-layer layout, an illustrative formula
is
\begin{equation}
 B_{\mathrm{physical}}=2bsd\sum_{\ell}H_{kv,\ell}\max_h n_{\ell h}.
\end{equation}
This formula is layout-specific; tuple caches or native paged layouts need
their own accounting. It explains why a route-local choice can fit in isolation
and still cause the complete cache transaction to exceed its cap.

Let $\mathcal P(C)$ be the set of physically feasible cache plans under cap $C$.
The shared-byte protocol asks a method to reach its largest attainable storage
level $B^*(C)=\max_{P\in\mathcal P(C)}B_{\mathrm{active}}(P)$, then optimize its
selection policy within that level. A greedy method's nonpositive marginal-gain
stopping rule is not an excuse to underfill this experimental control.
Mandatory, prompt, evidence, prototype-anchor, and per-head minimum reservations
must be reconciled in one feasibility transaction. Optional reservations can
be reduced deterministically; mandatory state and the hard cap cannot.

If two methods have different attainable lattices, the protocol either chooses
a common attainable target or reports the discrepancy and excludes the pair
from strict same-byte claims. It does not silently compare the nominal cap
while hiding different actual sizes. Native engine block boundaries are an
important source of such differences.

\subsection{Storage audit and atomic construction}
Before committing a plan, record knapsack-estimated bytes, unique selected-source
bytes, planner-tier bytes, external-exact bytes, constructed active bytes, and
physical padding bytes. These are related diagnostics, not six disjoint terms
to add. For a defined layout, constructed bytes must equal actual unique
active KV allocation; padding is the difference from populated active payload
under the same storage boundary.

Construct a new cache in isolation, verify its validity masks and logical
position map, then publish it to decode. Do not partially update a live cache
and leave scheduler or position state pointing at the old layout. A rejected
compressed proposal retains its accounting separately from any FullKV fallback.
A fallback has different memory and cannot satisfy the compression gate.

\subsection{Prototype construction}
For compatible source token vectors $k_j,v_j$ assigned weights
$\alpha_j\ge0$, $\sum_j\alpha_j=1$, a phase-neutral construction is
\begin{align}
 \bar k&=\sum_{j\in\mathcal C}\alpha_j R(p_j)^{-1}k_j,\\
 k_{\mathcal C}&=R(p_*)\bar k,\qquad
 v_{\mathcal C}=\sum_{j\in\mathcal C}\alpha_j v_j.
\end{align}
This expression applies only when $R$ is the adapter's verified invertible
RoPE transform and the key's cached representation is known. Partial rotary
dimensions, scaling variants, and multimodal coordinates need adapter-specific
handling. Non-rotated value vectors are averaged directly. Accumulation uses
FP32 before casting to the original cache dtype.

An anchor is a selected exact block used for assignment and position selection;
it is not silently mixed into the prototype and counted a second time. Merges
stay within a layer/head and enforce configured modality and positional-span
limits. Diagnostics record members, weights, modality composition, span,
key/value dispersion, and physical bytes saved after adding prototype slots.

\paragraph{Why phase safety is not equivalence.}
Original attention includes several exponentials,
$\sum_j\exp(q^\top k_j/\sqrt d)$, whereas a single prototype contributes one.
Averaging keys does not preserve that sum or its value-weighted numerator.
Even identical keys with multiple members need a multiplicity treatment to
preserve their mass relative to other tokens. MosaicKV therefore treats
prototype merging as a lossy approximation, never as a consequence-free
deduplication. Disable merging if representation compatibility cannot be
established. At retention one, return the unchanged full cache.

\subsection{Residual storage and recovery}
Residual storage maps layer, head, prototype ID, and original position to
original source blocks. The conservative format is lossless in the source
dtype and pinned on the CPU where supported. FP16, BF16, FP8, and INT8 variants
need explicit dtype/scale metadata and separate quality validation; lower byte
width is not automatically lower latency. Host pinning can consume scarce
locked memory and must have an error policy.

For next-token probabilities $p_v$, normalized entropy is
\begin{equation}
 H(p)=-\frac{\sum_{v=1}^{|V|}p_v\log p_v}{\log|V|}.
\end{equation}
A prototype-risk policy uses cluster mass $a_{\mathcal C}$ and dispersion
$D_{\mathcal C}$, with score
$\rho_{\mathcal C}=a_{\mathcal C}D_{\mathcal C}$. The manifest defines how key
and value dispersion combine, how heads/layers aggregate, and whether a cheap
draft KL is also used. An entropy-only policy need not inspect prototypes but
still requires available residual evidence to restore.

\begin{figure*}[t]
\centering\begin{tikzpicture}[node distance=7mm]
\node[box,text width=2.55cm] (exact) {\textbf{Exact}\\Original selected K/V\\GPU active storage};
\node[optbox,text width=2.55cm,right=of exact] (proto) {\textbf{Prototype}\\Weighted K/V\\GPU active storage};
\node[optbox,text width=2.55cm,right=of proto] (resid) {\textbf{Residual}\\Original members\\Pinned CPU storage};
\node[box,text width=2.55cm,right=of resid] (provisional) {Provisional token\\entropy / prototype risk};
\node[optbox,text width=2.55cm,below=13mm of provisional] (restore) {Rank and restore members\\wait for transfer};
\node[optbox,text width=2.55cm,left=of restore] (budget) {Promote exact blocks\\evict if budget requires};
\node[optbox,text width=2.55cm,left=of budget] (redo) {Rollback provisional state\\recompute once};
\node[box,text width=2.55cm,left=of redo] (commit) {Commit token and cache\\next logical position};
\draw[optional] (proto)--(resid);
\coordinate (decodebus) at ($(exact.north)+(0,0.6)$);
\draw[flow] (exact.north)--(decodebus)-| (provisional.north);
\draw[optional] (proto.north)--(proto.north|-decodebus);
\draw[optional] (provisional)--node[right,font=\scriptsize]{trigger}(restore);
\draw[optional] (restore)--(budget);
\draw[optional] (budget)--(redo);
\draw[optional] (redo)--(commit);
\draw[optional] (resid.south)--++(0,-0.4)-| (restore.west);
\draw[flow] (provisional.east)--++(0.4,0)|- ($(commit.south)+(0,-0.6)$)--(commit.south);
\node[font=\scriptsize,fill=white] at ($(redo.south)+(0,-0.6)$) {no trigger: commit provisional token};
\end{tikzpicture}

\caption{Compatibility-gated repair transaction. The current token is
provisional until transfer, budget enforcement, and at most one re-decode
complete. Dashed stages require separate validation.}
\label{fig:repair}
\end{figure*}

Asynchronous transfers use a stream/event dependency before restored blocks
are consumed. Promotion persists for later steps. If only some members of a
prototype are restored, its remaining membership and weights must be updated
or the prototype removed with explicit accounting; keeping a stale prototype
can double-count restored evidence. Non-promoted prototype eviction must not
remove mandatory or newly promoted exact states. Policy \texttt{none} executes
no repair. Each event records trigger reason, step, clusters, restored IDs and
bytes, transfer time, re-decode time, logit/token change, and evaluation-only
quality benefit.

\subsection{Bounded logit-safe probing as a separate variant}
The repository also contains exploratory hierarchical search and cache-isolation
machinery. A bounded probe can compare candidate next-token distributions to
FullKV using a fixed KL/logit threshold and greedy-token agreement. Route
ladders must contain only physically feasible levels. A byte allocator then
chooses route levels, followed by a final joint probe: individually safe changes
need not be safe together.

This is not a zero-cost verification oracle or a guarantee about the entire
answer. Probe count, horizon, rollback, failed candidate accounting, and
fallback status are part of the effective method. Repeated failed refinements
can make planning far more expensive than decode. Such a variant must be
preregistered and compared with the same byte and generation controls; it is
not silently substituted for the frozen primary candidate.
